\documentclass[letterpaper]{article}
\usepackage[margin=1in]{geometry}
\usepackage{amsmath,amssymb,amsthm}
\usepackage{natbib}
\usepackage{booktabs}
\usepackage{microtype}

\newcommand{\task}{\ensuremath{\Pi}}
\newcommand{\vars}{\ensuremath{\mathcal{V}}}
\newcommand{\states}{\ensuremath{\mathcal{S}}}
\newcommand{\sinit}{\ensuremath{s_{\mathrm{I}}}}
\newcommand{\cost}{\ensuremath{\mathit{cost}}}
\newcommand{\hstar}{\ensuremath{h^{*}}}
\newcommand{\hblind}{\ensuremath{h^{0}}}
\newcommand{\hpot}{\ensuremath{h^{\text{pot}}}}
\newcommand{\gstar}{\ensuremath{g^{*}}}
\newcommand{\copt}{\ensuremath{C^{*}}}
\newcommand{\bdd}[1]{\ensuremath{B_{#1}}}
\newcommand{\bddsize}[1]{\ensuremath{|B_{#1}|}}
\newcommand{\cofcount}{\ensuremath{Q}}
\newcommand{\width}{\ensuremath{\mathsf{w}}}
\newcommand{\values}{\ensuremath{V}}
\newcommand{\level}[2]{\ensuremath{H_{#2}}}
\newcommand{\effort}{\ensuremath{\mathit{effort}}}
\newcommand{\sub}[2]{\ensuremath{\mathrm{cof}_{#1}(#2)}}
\newcommand{\img}{\ensuremath{\mathrm{img}}}
\newcommand{\bddastar}{BDDA\ensuremath{^{*}}}
\newcommand{\astar}{A\ensuremath{^{*}}}
\newcommand{\ghseta}{GHSETA\ensuremath{^{*}}}

\theoremstyle{plain}
\newtheorem{theorem}{Theorem}
\newtheorem{lemma}{Lemma}
\newtheorem{proposition}{Proposition}
\newtheorem{corollary}{Corollary}
\theoremstyle{definition}
\newtheorem{definition}{Definition}
\newtheorem{example}{Example}
\theoremstyle{remark}
\newtheorem{remark}{Remark}

\title{Fixed-Order Residual Width Bounds for Forward
Product-at-Evaluation Symbolic \astar: Supplementary Material}
\author{Anonymous for review}
\date{}

\begin{document}
\maketitle

\section{Search Setup and Bucket Identity}

Let \task\ be a finite deterministic planning task with positive integer
operator costs, initial state \sinit, and optimal solution cost
$0<\copt<\infty$.
For a state \(s\), let \(\gstar(s)\) be its cheapest-path cost from
\sinit, let \(L_g=\{s\mid\gstar(s)=g\}\), and let
\(\level{h}{v}=\{s\mid h(s)=v\}\).  The heuristic
\(h\colon\mathcal{S}\to\mathbb{Z}_{\geq0}\cup\{\infty\}\) is consistent if
\(h(s)\leq\cost(o)+h(s[o])\) for every applicable operator \(o\), and
\(h=0\) on goal states.  In \bddastar, states with infinite heuristic
value are discarded; the remaining generated states are grouped by keys
\((g,v)\), selected by increasing \(f=g+v\) and then increasing \(g\).
Closed states are removed, and search terminates when a selected bucket
contains a goal.  This is the setting of the following restatement.

\begin{lemma}[bucket identity]
\label{lem-identity}
Let $h$ be consistent. \bddastar\ on \task\ with $h$ selects every
bucket key $(g,v)$ with $g + v \leq \copt$ and $g < \copt$ and
$L_g \cap \level{h}{v} \neq \emptyset$ exactly once before terminating,
and the BDD expanded with key $(g,v)$ represents exactly
$L_g \cap \level{h}{v}$. No other bucket is expanded before
termination.
\end{lemma}

\begin{proof}
Throughout, the \emph{key} of a generated state $s$ is the pair
$(c, h(s))$, where $c$ is the cost of the generating path, and the
$f$-value of a key $(g,v)$ is $g + v$. We prove five claims.

\emph{Claim 1 (monotone selection): if a key $(g_1,v_1)$ is selected at
step $t_1$ and a key $(g_2,v_2)$ at step $t_2 > t_1$, then
$g_1 + v_1 \leq g_2 + v_2$, and if $g_1 + v_1 = g_2 + v_2$ then
$g_1 \leq g_2$.}
For the $f$-part, every state inserted into the open list during the
expansion of a bucket with key $(g,v)$ is a successor $s[o]$ of a state
$s$ with $h(s) = v$ and has key $(g + \cost(o),\, h(s[o]))$, whose
$f$-value satisfies
$g + \cost(o) + h(s[o]) \geq g + h(s) = g + v$ by consistency. Hence
expanding a minimum-$f$ key only inserts keys of at least that
$f$-value, the minimum $f$-value of the open list never decreases, and
selected $f$-values are non-decreasing over time. For the $g$-part,
suppose for contradiction that there are selections at steps
$t_1 < t_2$ of keys with equal $f$-value $F$ and $g_2 < g_1$, and among
all such violations choose one minimizing $t_2$. The bucket of
$(g_2,v_2)$ must be empty at step $t_1$: otherwise the tie-breaking
rule would select it instead of $(g_1,v_1)$. Its content is also not
inserted at step $t_1$ itself, since states inserted by the expansion
of $(g_1,v_1)$ have $g$-component $g_1 + \cost(o) > g_1 > g_2$. So
every state of the bucket is inserted at some step $t'$ with
$t_1 < t' < t_2$ by expanding a key $(g',v')$ with
$g' = g_2 - \cost(o) < g_2$ for some operator $o$. Consistency gives
$g' + v' \leq g_2 + v_2 = F$ as above, and the $f$-part applied to
$t_1 < t'$ gives $g' + v' \geq F$, so $g' + v' = F$ and $g' < g_1$:
the pair $(t_1, t')$ is a violation with $t' < t_2$, contradicting
minimality.

\emph{Claim 2 (no re-population): after a key $(g,v)$ is selected at
step $t$, no state with key $(g,v)$ is inserted into the open list;
consequently every key is selected at most once, and at its selection
its bucket contains every state that was ever generated with this key
and not expanded before.}
Suppose a state with key $(g,v)$ is inserted at a step $t' > t$ by
expanding a key $(g',v')$. Then $g' + v' \leq g + v$ by consistency as
in Claim~1, and $g' + v' \geq g + v$ by the $f$-part of Claim~1, so the
$f$-values are equal; but $g' = g - \cost(o) < g$ contradicts the
$g$-part. A second selection of the same key would require an
insertion after step $t$, which is excluded. Finally, a state generated
with key $(g,v)$ before step $t$ enters the bucket unless it is removed
by the closed-list test, which removes only previously expanded states,
and states are removed from the open list only when their bucket is
selected.

\emph{Claim 3 (optimal $g$-values): every state $s$ in a selected
bucket with key $(g,v)$ satisfies $\gstar(s) = g$ and $h(s) = v$.}
We use induction on the selection step; Claims~1 and~2 hold
unconditionally and may be used freely. That $h(s) = v$ holds by
construction of the level-set partition, and $\gstar(s) \leq g$ because
$s$ is reached by a path of cost $g$. Suppose $\gstar(s) < g$ for some
$s$ in the bucket selected at the current step, and fix a cheapest path
$p_0 = \sinit, \dots, p_k = s$, all of whose prefixes are cheapest
paths to their end states. Let $j$ be the smallest index such that
$p_j$ has not been expanded before the current step; $j$ exists because
$s$ qualifies, and $j \geq 1$ because \bddastar\ selects the initial
bucket $\{\sinit\}$ first. By the induction hypothesis, every earlier
expansion of $p_{j-1}$ has key $(\gstar(p_{j-1}), h(p_{j-1}))$, so
$p_j$ was generated with key $(\gstar(p_j), h(p_j))$ before the current
step. This key has not been selected yet: a selection after the
generation of $p_j$ would have expanded $p_j$ (Claim~2), contradicting
the choice of $j$, and a selection before the generation would
contradict Claim~2 directly. Hence the open list currently contains the
key $(\gstar(p_j), h(p_j))$, and consistency along the path segment
from $p_j$ to $s$ yields
$\gstar(p_j) + h(p_j) \leq \gstar(s) + h(s) < g + v$, using
$\gstar(s) < g$ and $h(s) = v$. This contradicts minimum-$f$ selection.

\emph{Claim 4 (completeness): every state $s$ with
$\gstar(s) = g < \copt$ and $g + h(s) \leq \copt$ is generated with key
$(g, h(s))$, and this key is selected before the search terminates.}
By strong induction on $g$. \emph{Generation:} for $g = 0$, $s = \sinit$
initializes the open list with key $(0, h(\sinit))$. For $g > 0$, let
$p$ be the predecessor of $s$ on a cheapest path, with
$\gstar(p) = g - \cost(o) < g$ and, by consistency,
$\gstar(p) + h(p) \leq \gstar(s) + h(s) \leq \copt$. By the induction
hypothesis, $p$ is generated with key $(\gstar(p), h(p))$ and this key
is selected before termination; by Claim~2 the bucket then contains
$p$, because any earlier expansion of $p$ would, by Claim~3, have
occurred with the same key, which is selected only once. Expanding the
bucket generates $s$ at cost $\gstar(p) + \cost(o) = g$ with key
$(g, h(s))$. \emph{Selection:} let $K = (g, h(s))$, with
$f(K) \leq \copt$ and $g < \copt$. From the step at which $K$ first
becomes nonempty, $K$ remains in the open list until it is selected
(Claim~2). The search terminates only when a selected bucket contains a
goal state, and by Claim~3 a goal state in a selected bucket has key
$(\gstar(s_G), 0)$ with $\gstar(s_G) \geq \copt$, since consistent
heuristics assign $0$ to goal states and \copt\ is the cheapest goal
cost. While $K$ is open, no such key is selected: its $f$-value is at
least $\copt \geq f(K)$, and in case of equality its $g$-component
$\gstar(s_G) \geq \copt > g$ loses the tie-break against $K$ or another
open key with smaller $g$-component. Hence the search does not
terminate while $K$ is open. Moreover, while $K$ is open every
selection picks a key with $f$-value at most $f(K) \leq \copt$; there
are at most $(\copt + 1)(\copt + 2)/2$ keys $(g', v')$ of nonnegative
integers with $g' + v' \leq \copt$, each selected at most once
(Claim~2), and the open list is nonempty while $K$ is open. Hence after
finitely many steps $K$ itself is selected, before termination.

\emph{Claim 5 (no other expansions): every bucket selected before the
terminating selection has a key $(g,v)$ with $g + v \leq \copt$ and
$g < \copt$.}
Fix a cheapest plan $q_0 = \sinit, \dots, q_r$ with $q_r$ a goal state
and all prefixes cheapest, so that consistency yields
$\gstar(q_l) + h(q_l) \leq \copt$ for every $l$. First we exhibit a
witness key: at every step at which the search has not yet terminated,
let $l$ be the smallest index such that $q_l$ is unexpanded, which
exists because $q_r$ is unexpanded before termination. If $l = 0$, the
initial key $(0, h(\sinit))$ has not been selected, since its bucket
contains $\sinit$, and is therefore open. If $l \geq 1$, then $q_{l-1}$
was expanded, by Claim~3 with key $(\gstar(q_{l-1}), h(q_{l-1}))$, so
$q_l$ was generated with key $K_l = (\gstar(q_l), h(q_l))$; exactly as
in the proof of Claim~3, $K_l$ has not been selected and is therefore
open. In all cases the open list contains a key with $f$-value at most
\copt, namely $(0, h(\sinit))$ or $K_l$; call it the witness. Second,
minimum-$f$ selection against the witness implies that every selected
key satisfies $g + v \leq \copt$. Third, a selected key with
$g + v \leq \copt$ and $g \geq \copt$ must equal $(\copt, 0)$, and we
show that its selection is the terminating one. Consider the step at
which $(\copt, 0)$ is selected and let $l$ be as above. If $l < r$,
the witness $K_l$ is open with $f$-value at most \copt\ and
$g$-component $\gstar(q_l) < \copt$; then either its $f$-value is
smaller than \copt, contradicting minimum-$f$ selection of
$(\copt, 0)$, or it equals \copt\ and the tie-breaking rule prefers
$K_l$ or another open key of smaller $g$-component over $(\copt, 0)$,
again a contradiction. Hence $l = r$, so $q_r$ was generated with key
$(\copt, 0)$ and, being unexpanded, lies in its bucket at selection
time by Claim~2. The selected bucket thus contains the goal state
$q_r$ and the search terminates. Consequently, every selection before
the terminating one has $g + v \leq \copt$ and $g < \copt$.

\emph{Assembly.} Let $(g,v)$ with $g + v \leq \copt$ and $g < \copt$
and $L_g \cap \level{h}{v} \neq \emptyset$. Every
$s \in L_g \cap \level{h}{v}$ satisfies the premise of Claim~4, so its
key $(g,v)$ is selected exactly once before termination (Claims~2
and~4), and at this selection the bucket contains $s$ (Claim~2, with
prior expansion of $s$ excluded because, by Claim~3, it would have
occurred at the same, uniquely selected key). Conversely, by Claim~3
the selected bucket contains only states with $\gstar = g$ and
$h$-value $v$. Hence the expanded BDD represents exactly
$L_g \cap \level{h}{v}$. Keys whose intersection is empty are never
selected, since buckets enter the open list only when populated, and by
Claim~5 no key outside the stated range is expanded before the search
stops.
\end{proof}

\section{Fixed-Order Residual-Width Notation and Auxiliary Results}

The maximum number of distinct residual functions at a fixed-order cut is the
standard width measure for ordered branching programs and MTBDDs
\citep{wegener-2000}. We use
\emph{cofactor width} as a planning-facing name for this existing measure.
The novelty claims concern its application to heuristic fragmentation and
forward product-at-evaluation search, the separation constructions, and the
planning-heuristic family bounds, not the width definition.

Fix Boolean variables $x_1,\dots,x_n$ in an order $\pi$.  For a Boolean
or integer-valued function $f$ and a prefix assignment
$\rho\in\{0,1\}^i$, write $f|_\rho$ for the residual function after the
prefix and
$\sub{i}{f}=\{f|_\rho\mid\rho\in\{0,1\}^i\}$.  The cofactor count and
cofactor width are
\[
  \cofcount(f)=\sum_{i=0}^{n-1}|\sub{i}{f}|,
  \qquad
  \width_\pi(f)=\max_{0\leq i\leq n}|\sub{i}{f}|.
\]
For a Boolean function $f$, $\bdd{f}$ is its reduced ordered BDD under
$\pi$, and $\bddsize{f}$ counts its inner nodes.  For a heuristic $h$,
$\img(h)$ is its attained value set.  The blind heuristic $\hblind$ is
identically zero, and $\effort(\task,h)$ is the cumulative number of inner
nodes in the BDDs expanded by product-at-evaluation forward \bddastar\ before
optimal termination.
Each SAS$^+$ variable $v$ is encoded using
$\lceil\log_2|D_v|\rceil$ Boolean bits; the order $\pi$ is
variable-contiguous when all bits of each variable are consecutive. This
encoding may leave invalid Boolean assignments. A heuristic on states is
equipped with a fixed total extension to the full Boolean cube; where
relevant below, the statements specify that extension.

\begin{proposition}[value range]
\label{prop-range}
If $\width_\pi(h) \leq W$, then $|\img(h)| \leq W$.
\end{proposition}

\begin{proof}
The elements of $\sub{n}{h}$ are constant functions, one for each attained
value. Hence $|\img(h)|=|\sub{n}{h}|\leq\width_\pi(h)$.
\end{proof}

\begin{proposition}[closure]
\label{prop-closure}
Let $h_1,h_2$ have widths $W_1,W_2$ under $\pi$ and let
$\oplus\in\{+,\max,\min\}$. Then
$\width_\pi(h_1\oplus h_2)\leq W_1W_2$.
\end{proposition}

\begin{proof}
At cut $i$, each cofactor of $h_1\oplus h_2$ is the pointwise combination
of a realized pair $(h_1|_\rho,h_2|_\rho)$. There are at most
$|\sub{i}{h_1}|\,|\sub{i}{h_2}|\leq W_1W_2$ such pairs. Maximizing over
$i$ proves the claim.
\end{proof}

Call $q\colon\states\to\mathbb Z\cup\{\infty\}$ a \emph{raw admissible
and consistent estimate} if $q\leq\hstar$, if
$q(s)\leq c+q(s')$ for every transition $s\xrightarrow{c}s'$, and if
$q\leq0$ on goals.

\begin{proposition}[safe terminal transforms]
\label{prop-transform}
Let $q$ be a raw admissible and consistent estimate, equipped with a total
extension $q\colon\{0,1\}^n\to\mathbb Z\cup\{\infty\}$. Let
$\phi\colon\mathbb Z\cup\{\infty\}\to
\mathbb Z_{\geq0}\cup\{\infty\}$ be nondecreasing, satisfy
$\phi(0)=0$, $\phi(\infty)=\infty$, and
\[
  \phi(a+c)\leq\phi(a)+c
  \quad\text{for every $a\in\mathbb Z$ and integer $c\geq0$.}
\]
Then $h=\phi\circ q$ is admissible and consistent and
$\width_\pi(h)\leq\width_\pi(q)$.
\end{proposition}

\begin{proof}
For a transition $s\xrightarrow{c}s'$ with finite $q(s')$,
\[
  \phi(q(s))\leq\phi(c+q(s'))\leq c+\phi(q(s'));
\]
if $q(s')=\infty$, consistency is immediate. Monotonicity,
nonnegativity, and $\phi(0)=0$ give $\phi(a)=0$ for $a\leq0$, while the
cost condition with $a=0$ gives $\phi(a)\leq a$ for $a\geq0$. Thus, for
finite $q(s)$,
$\phi(q(s))\leq\max(0,q(s))\leq\hstar(s)$. If $q(s)=\infty$, raw
admissibility implies $\hstar(s)=\infty$. The transformed heuristic is
therefore admissible, consistent, and zero on goals. Finally,
$(\phi\circ q)|_\rho=\phi\circ(q|_\rho)$, so the terminal map may merge
cofactors but cannot create any.
\end{proof}

Rectification $\phi_+(a)=\max(0,a)$ is one instance. For
$\kappa\in\mathbb Z_{\geq0}$, define
\[
\phi_\kappa(a)=
\begin{cases}
\min(\max(0,a),\kappa),&a<\infty,\\
\infty,&a=\infty.
\end{cases}
\]
It satisfies the proposition and has at most $\kappa+1$ finite values.
Moreover, it is pointwise strongest among transforms satisfying the
proposition whose finite outputs lie in $\{0,\ldots,\kappa\}$: the proof
above and the output ceiling give
$\phi(a)\leq\min(\max(0,a),\kappa)=\phi_\kappa(a)$.

Let $\mathrm{ADD}_\pi(h)$ be the reduced ordered ADD of $h$ under $\pi$
\citep{bahar-et-al-iccad1993}, and let $|\mathrm{ADD}_\pi(h)|$ count
inner nodes.

\begin{proposition}[ADD characterization]
\label{prop-add}
$\width_\pi(h)\leq|\mathrm{ADD}_\pi(h)|+|\img(h)|$ and
$|\mathrm{ADD}_\pi(h)|\leq n\width_\pi(h)$.
\end{proposition}

\begin{proof}
Every cofactor at a cut is either a constant or the function computed by a
unique inner ADD node below the cut, giving the first inequality. Every
inner node labeled $x_{i+1}$ computes a distinct element of
$\sub{i}{h}$, so there are at most $\width_\pi(h)$ nodes at each of the
$n$ levels, giving the second.
\end{proof}

\begin{remark}[total extensions]
\label{rem-encoding}
Cofactor width is measured on the chosen total extension, while a lower
bound witnessed entirely by genuine states applies to every such extension.
\end{remark}

\section{Reduced and Quasi-Reduced OBDD Size}

\begin{lemma}
\label{lem-qsize}
For every Boolean $f$ over $x_1, \dots, x_n$,
$\bddsize{f} \leq \cofcount(f) \leq n \cdot (\bddsize{f} + 2)$.
\end{lemma}

\begin{proof}
For the first inequality, consider a node $N$ of $\bdd{f}$ labeled
$x_{i+1}$. Since every node of a reduced OBDD is reachable from the root,
there is a path from the root to $N$; fixing the variables tested along
this path to the corresponding edge values and all remaining variables
among $x_1, \dots, x_i$ arbitrarily yields a prefix $\rho \in \{0,1\}^i$
with $f|_\rho$ equal to the function computed by $N$, because variables
skipped by long edges do not influence the traversal. Hence every node
labeled $x_{i+1}$ computes an element of $\sub{i}{f}$, and by canonicity
distinct nodes compute distinct functions, so the map is injective.
Summing over $i$ gives
$\bddsize{f} \leq \sum_{i=0}^{n-1} |\sub{i}{f}| = \cofcount(f)$. For the
second inequality, fix $i$ and consider any $u \in \sub{i}{f}$. Either
$u$ is a constant function, of which there are at most two, or $u$
essentially depends on some variable; let $x_j$ with $j > i$ be the first
such variable. Traversing $\bdd{f}$ from the root under any prefix
witnessing $u$ ends at an inner node labeled $x_j$ that computes $u$, and
by canonicity this node is unique. Hence
$|\sub{i}{f}| \leq \bddsize{f} + 2$ for every $i$, and summing over the
$n$ levels yields the claim.
\end{proof}

\section{Fragmentation Bounds}

Fix a heuristic $h$ of width $W$ under $\pi$ and let
$T=|\img(h)|$.

\begin{lemma}[pair identity]
\label{lem-pair}
Let $S\subseteq\{0,1\}^n$, let $U\subseteq\img(h)$, and let
$\rho,\rho'\in\{0,1\}^i$ satisfy
$\chi_S|_\rho=\chi_S|_{\rho'}$ and $h|_\rho=h|_{\rho'}$. Then
$\chi_{S\cap h^{-1}(U)}|_\rho=
\chi_{S\cap h^{-1}(U)}|_{\rho'}$.
\end{lemma}

\begin{proof}
For every suffix $\sigma\in\{0,1\}^{n-i}$,
\[
 \chi_{S\cap h^{-1}(U)}|_\rho(\sigma)
 =\chi_S|_\rho(\sigma)[h|_\rho(\sigma)\in U].
\]
Both factors are unchanged when $\rho$ is replaced by $\rho'$, proving
equality of the residual functions.
\end{proof}

\begin{theorem}[bucket bound]
\label{thm-bucket}
Let $S\subseteq\{0,1\}^n$ and $U\subseteq\img(h)$. Then
\[
\bddsize{S\cap h^{-1}(U)}
\leq\cofcount(\chi_{S\cap h^{-1}(U)})
\leq W\cofcount(\chi_S)
\leq Wn(\bddsize{S}+2).
\]
\end{theorem}

\begin{proof}
The first and last inequalities follow from Lemma~\ref{lem-qsize}. Fix a
cut $i$ and let
\[
 R_i=\{(\chi_S|_\rho,h|_\rho)\mid\rho\in\{0,1\}^i\}.
\]
For each $u\in\sub{i}{\chi_{S\cap h^{-1}(U)}}$, choose a witnessing
prefix $\rho_u$ and map $u$ to
$(\chi_S|_{\rho_u},h|_{\rho_u})\in R_i$. This map is injective by
Lemma~\ref{lem-pair}. Consequently,
\[
 |\sub{i}{\chi_{S\cap h^{-1}(U)}}|
 \leq|R_i|
 \leq|\sub{i}{\chi_S}|\,|\sub{i}{h}|
 \leq W|\sub{i}{\chi_S}|.
\]
Summing over $i=0,\ldots,n-1$ proves the middle inequality.
\end{proof}

\begin{theorem}[partition bound]
\label{thm-partition}
Let $S\subseteq\{0,1\}^n$. Then
\[
 \sum_{v\in\img(h)}\bddsize{S\cap\level{h}{v}}
 \leq TW\cofcount(\chi_S)
 \leq TWn(\bddsize{S}+2)
 \leq W^2n(\bddsize{S}+2).
\]
If $S\subseteq\states$ and the sum ranges only over finite values attained
on planning states, the first two bounds hold with $T$ replaced by
$\values=|h(\states)\setminus\{\infty\}|$.
\end{theorem}

\begin{proof}
Theorem~\ref{thm-bucket} with $U=\{v\}$ bounds every bucket by
$W\cofcount(\chi_S)$. Sum over the $T$ terminals and use
$T\leq W$ from Proposition~\ref{prop-range}; Lemma~\ref{lem-qsize}
converts cofactor count to reduced size. If only finite values attained on
planning states are included, there are at most $\values$ nonempty
buckets.
\end{proof}

\begin{example}[the factor $n$ is necessary]
\label{ex-parity}
For $n\geq2$, let $S=\{s\mid s(x_1)=1\}$ and
$h(s)=(\sum_{i=1}^ns(x_i))\bmod2$. Then $\bddsize{S}=1$ and $h$ has
width two under any order. The bucket $S\cap\level{h}{0}$ fixes
$x_1=1$ and requires odd parity on $x_2,\ldots,x_n$. Its reduced BDD has
one node for $x_1$, one for $x_2$, and two for each of
$x_3,\ldots,x_n$, hence $2n-2$ inner nodes. Thus a width-two slice can be
$\Theta(n)$ times larger than its source set. In cofactor count,
$\cofcount(\chi_S)=2n-1$, so Theorem~\ref{thm-bucket} is within a
constant factor.
\end{example}

\section{Search-Effort Bound}

By Lemma~\ref{lem-identity}, expansion size is exactly
\[
 \effort(\task,h)
 =\sum_{g<\copt}\ \sum_{\substack{v\in\img(h)\\g+v\leq\copt}}
   \bddsize{L_g\cap\level{h}{v}}.
\]

\begin{theorem}[effort bound]
\label{thm-effort}
Let $h$ be consistent with $\width_\pi(h)\leq W$, and let
$\values=|h(\states)\setminus\{\infty\}|$. Then
\[
 \effort(\task,h)
 \leq\values Wn\bigl(\effort(\task,\hblind)+2\copt\bigr)
 \leq W^2n\bigl(\effort(\task,\hblind)+2\copt\bigr).
\]
\end{theorem}

\begin{proof}
The bucket identity permits enlarging the inner sum to all finite values
attained on planning states. The finite-state form of
Theorem~\ref{thm-partition} therefore gives
\[
 \effort(\task,h)
 \leq\sum_{\substack{g<\copt\\L_g\ne\emptyset}}
       \values Wn(\bddsize{L_g}+2).
\]
Let $N$ be the number of nonempty layers below $\copt$. Since blind
search expands each such layer once,
\begin{equation}
 \effort(\task,h)
 \leq\values Wn\bigl(\effort(\task,\hblind)+2N\bigr).
 \label{eq-effort-nonempty}
\end{equation}
Positive integer costs imply $N\leq\copt$: the possible layer indices are
$0,\ldots,\copt-1$. This proves the first inequality, and
$\values\leq W$ proves the second.
\end{proof}

\begin{corollary}[search-effort ratio]
\label{cor-effort}
Under the assumptions of Theorem~\ref{thm-effort},
\begin{equation}
 \effort(\task,h)
 \leq3\values Wn\,\effort(\task,\hblind)
 \leq3W^2n\,\effort(\task,\hblind).
 \label{eq-supp-effort-ratio}
\end{equation}
\end{corollary}

\begin{proof}
Positive costs give $L_0=\{\sinit\}$. Because
$0<\copt<\infty$, this is a proper subset of the Boolean cube. Every
nonempty $L_g$ with $0<g<\copt$ is also proper because it omits
\sinit. Its characteristic function is therefore nonconstant and its
reduced OBDD has at least one inner node. Hence
$N\leq\effort(\task,\hblind)$ in
Equation~\eqref{eq-effort-nonempty}, which proves both inequalities.
\end{proof}

This is a one-sided bound on expanded frontier representations in forward
product-at-evaluation \bddastar. It neither bounds the primitive relational
products of delta-partitioned methods nor gives a cumulative bidirectional
guarantee; it does allow unbounded savings through $f$-cutoff and dead-end
pruning.

\section{Image Calls and Pruning-Only Search}

\begin{definition}[image-aware effort]
For $\alpha\geq0$, the \emph{$\alpha$-effort} of a symbolic run is
\[
 \effort_\alpha=\sum_B(\alpha+|B|),
\]
where the sum ranges over idealized high-level cost-image invocations before
termination and $B$ is the argument of the invocation.
\end{definition}

For $\alpha=0$, this is the expansion-size measure when the high-level
cost-image routine is invoked once for each expanded bucket in this model.

\begin{proposition}[image-call counts]
\label{prop-calls}
Let $h$ be consistent with $\width_\pi(h)\leq W$, and let $N_\ell$ be
the number of nonempty semantic blind layers $L_g$ with $g<\copt$. An
idealized blind implementation storing each semantic $L_g$ in one BDD performs
$N_\ell$ counted high-level cost-image invocations, one per nonempty semantic
layer. An idealized product-at-evaluation
\bddastar\ performing one high-level cost-image invocation per expanded
nonempty $(g,v)$ bucket performs $\sum_g\values_g$ counted invocations, where
$\values_g$ is the number of nonempty expanded buckets in layer $g$ and
$\values_g\leq\values\leq W$; hence it performs at most
$\values N_\ell\leq WN_\ell$ counted invocations.
\end{proposition}

\begin{proof}
The idealized blind implementation invokes the high-level cost-image routine
once for each nonempty semantic layer.
Lemma~\ref{lem-identity} makes every active heuristic layer a blind layer
below $\copt$, and partitions it into exactly $\values_g$ nonempty
buckets. Proposition~\ref{prop-range} gives
$\values_g\leq\values\leq W$.
\end{proof}

Thus the invocation-overhead term is at most $\alpha WN_\ell$, but width does
not provide a width-independent one-call-per-layer guarantee. Here one
count means one idealized invocation of the high-level cost-image routine per
semantic layer or bucket, not every relational product it performs. In
particular, the proposition does not
apply verbatim to SetA$^*$ or cost-general \ghseta: operator-potential
implementations partition transition relations by operator-cost and
heuristic-delta pairs, so their relational-product count also depends on
those partitions. The exact finite value count $\values$ can therefore be
a much sharper operational parameter than an ADD-derived width bound, while
delta-integrated search requires a separate analysis
\citep{fiser-et-al-aij2024}. On the quadratic construction below, the
idealized model counts $\Theta(W)$ high-level cost-image invocations on each
of $\Theta(k)$ late semantic layers, so the factor is realizable.

\emph{Pruning-only search} instead maintains one BDD per uniform-cost
layer and an incumbent solution bound $c$ ($c=\infty$ initially). It
discards $h=\infty$ and states satisfying $g+h\geq c$, and reapplies the
current bound when a layer is expanded. It stops when the next layer cannot
improve the incumbent.

\begin{lemma}[slice identity]
\label{lem-slice}
Let $h$ be consistent, and let $c_g$ be the incumbent when pruning-only
search expands layer $g$. The expanded layer is exactly
\[
 L_g\cap h^{-1}(\{v\mid v<c_g-g\}).
\]
\end{lemma}

\begin{proof}
Applying $c_g$ at expansion gives containment in the displayed slice. For
the converse, let $s\in L_g$ satisfy $g+h(s)<c_g$, and take a cheapest
path $p_0=\sinit,\ldots,p_k=s$. Consistency gives
$\gstar(p_j)+h(p_j)\leq g+h(s)<c_g$ for every prefix state. Its layer was
expanded no later than $g$, when the nonincreasing incumbent was at least
$c_g$, so $p_j$ survived generation-time and expansion-time pruning.
Induction along the path shows that $s$ is generated at cost $g$ and
survives into the expanded layer.
\end{proof}

\begin{corollary}[pruning bound]
\label{cor-prune}
Let $h$ be admissible and consistent with $\width_\pi(h)\leq W$. The
cumulative size of the layers expanded by pruning-only search is at most
\[
 Wn\bigl(\effort(\task,\hblind)+2\copt\bigr),
\]
and it performs one idealized high-level cost-image invocation per expanded
semantic layer. With $c=\infty$
throughout and no recognized dead ends, its semantic layers and expansions
coincide with blind forward search, although heuristic construction and
evaluation can still consume time.
\end{corollary}

\begin{proof}
By Lemma~\ref{lem-slice}, every expanded BDD is a single threshold slice
of $L_g$. Theorem~\ref{thm-bucket} bounds it by
$Wn(\bddsize{L_g}+2)$. Summing over the at most $\copt$ nonempty layers
below the optimal goal cost proves the node bound. Because no layer is
partitioned, exactly one modeled high-level cost-image invocation is made for
each expanded semantic layer. If the incumbent stays infinite and no state
has $h=\infty$, the retained slice
is all of $L_g$, proving the final claim.
\end{proof}

For every $\alpha\geq0$, pruning-only $\alpha$-effort is therefore at
most
\[
 Wn\bigl(\effort(\task,\hblind)+2\copt\bigr)+\alpha N_\ell.
\]
The discard is also pointwise sound in bidirectional blind search if the
forward direction uses an admissible goal-distance estimate and the
backward direction an admissible distance-from-initial estimate. We do not
lift this pointwise fact to a cumulative bidirectional bound: pruning
changes both closed sets and their meeting points, which can alter solution
discovery, incumbent updates, and an adaptive direction schedule. Those
feedback effects prevent identification with fixed blind reference
frontiers. A single linear merge-and-shrink abstraction can supply both
goal-distance and distance-from-initial estimates from the same shrunk
transition system. Pointwise incumbent pruning itself uses only
admissibility and, unlike the cumulative layer argument above, does not
require positive operator costs.

\section{Quadratic Width-Ratio Construction}

\begin{theorem}[quadratic width-ratio family]
\label{thm-tight}
For all integers $W \geq 2$ and $k \geq 4W$ there are tasks
$\task_{k,W}$ with $n = \Theta(k)$ Boolean variables and consistent,
admissible heuristics $h$ with $\width_\pi(h) = W$ such that
$\effort(\task_{k,W}, h) = \Omega(W^{2}) \cdot
\effort(\task_{k,W}, \hblind)$.
\end{theorem}

\begin{proof}
The task has binary variables $x_0,\dots,x_k$, choice variables
$v_1,\dots,v_k$, and one goal bit $t$. Initially only $x_0$ holds. For
each $i\in[k]$, two unit-cost operators advance the chain:
$o_i^0$ has precondition $x_{i-1}\wedge\neg x_i$ and effect $x_i$,
while $o_i^1$ has the same precondition and effect
$x_i\wedge v_i$. The only operator that achieves $t$ has precondition
$x_k\wedge\neg t$, effect $t$, and cost $W$. The goal is $t$, so every
solution has cost
\[
\copt=k+W.
\]
The order $\pi$ lists the $x$-block, then $v_1,\dots,v_k$, then $t$.

Let
\[
r(s)=\Bigl(\sum_{j=1}^{k}s(v_j)\Bigr)\bmod W,
\qquad
q(r)=(-r)\bmod W\in\{0,\dots,W-1\},
\]
and define
\[
h(s)=
\begin{cases}
0, & s(t)=1,\\
q(r(s)), & s(t)=0.
\end{cases}
\]
A unit-cost choice operator either leaves the residue unchanged or maps
$r$ to $r+1\bmod W$, and
$q(r)\leq1+q(r+1\bmod W)$. The goal operator can decrease $h$ by at
most $W-1$ and costs $W$. Hence $h$ is consistent. Every plan from a
non-goal state must still apply the cost-$W$ goal operator. Therefore a
solvable non-goal state $s$ satisfies
$h(s)\leq W-1<W\leq\hstar(s)$; admissibility is immediate at goals and
dead ends.

The heuristic ignores the $x$-block. While the $v$-block is read, its
cofactor is determined by the partial residue, of which there are at
most $W$; after $t$ is read the cofactors are constants. All $W$ values
are attained, so Proposition~\ref{prop-range} gives the reverse
inequality and
\[
\width_\pi(h)=W.
\]

\emph{Blind effort.}
For $i\leq k$, layer $L_i$ fixes the chain prefix, fixes
$v_j=0$ for $j>i$ and $t=0$, and leaves $v_1,\dots,v_i$ free. Therefore
$\bddsize{L_i}=\Theta(k)$. There are $k+1$ such pre-goal layers and no
intermediate layers between costs $k$ and $k+W$, so
\[
\effort(\task_{k,W},\hblind)=\Theta(k^2).
\]

\emph{Heuristic effort.}
On $L_i$, $h=q(\sum_{j\leq i}v_j\bmod W)$. For $i\geq W-1$ all $W$
buckets are nonempty, and every bucket is selected because
$i+h\leq k+W-1<\copt$. Fix a layer
$i\in[\lceil3k/4\rceil,k]$ and one bucket. For every
$W\leq j\leq i-W$, prefixes of $v_1,\dots,v_j$ with distinct partial
residues induce distinct cofactors of the modular bucket: a suffix using
at most $W-1$ of the remaining variables completes one residue to the
required bucket and not another. Each such cofactor essentially depends
on $v_{j+1}$ by the same argument. Consequently the reduced OBDD has
$W$ distinct nodes at each of the $i-2W+1=\Omega(k)$ middle levels, and
the bucket has size $\Omega(kW)$. Since $k\geq4W$, the interval
$[\lceil3k/4\rceil,k]$ contains $\Theta(k)$ layers. Summing over these
layers and their $W$ buckets gives
\[
\effort(\task_{k,W},h)
 \geq \Theta(k)\cdot W\cdot\Omega(kW)
 =\Omega(k^2W^2).
\]
Together with the blind bound, this proves the claimed
$\Omega(W^2)$ ratio.
\end{proof}

\section{Width of the Known Exponential Pathology}

For completeness, the family $\task_n$ has binary variables
$v_1,\dots,v_{2n}$ and $x_0,\dots,x_{2n+1}$. Initially only $x_0$ is
true. Operators $o_i$ and $\bar o_i$, whose precondition is
$x_{i-1}\wedge\neg x_i$, set $x_i$ true while setting $v_i$ true or
false, respectively. Once $x_0,\dots,x_{2n}$ are true, a goal operator
$g_i$ is applicable when $v_i\wedge v_{n+i}$ holds and achieves
$x_{2n+1}$. The order is
$\pi_n=x_0,\dots,x_{2n+1},v_1,\dots,v_{2n}$.

\begin{theorem}[width lower bound]
\label{thm-lb}
For the tasks $\task_n$ and orders $\pi_n$ of
\citet{speck-et-al-icaps2020},
$\width_{\pi_n}(\hstar) \geq 2^{n}$.
\end{theorem}

\begin{proof}
We exhibit $2^n$ prefixes with pairwise distinct \hstar-cofactors. For
$A \subseteq \{1, \dots, n\}$ let $\rho_A$ be the prefix that assigns
the $x$-variables the encoding of chain position $2n$, i.e.,
$x_0, \dots, x_{2n}$ true and $x_{2n+1}$ false, and assigns $v_i = 1$
iff $i \in A$ for $i \in \{1, \dots, n\}$. All $\rho_A$ fix the same
variables $x_0, \dots, x_{2n+1}, v_1, \dots, v_n$, so their cofactors
lie in $\sub{m}{\hstar}$ for the same level $m = (2n + 2) + n$. Let
$A \neq A'$ and pick without loss of generality
$i \in A \setminus A'$. Define the suffix $\sigma$ over
$v_{n+1}, \dots, v_{2n}$ by $v_{n+i} = 1$ and $v_{n+j} = 0$ for
$j \neq i$. The state $\rho_A \sigma$ satisfies
$v_i \wedge v_{n+i}$, so a goal operator is applicable and
$\hstar(\rho_A \sigma) = 1$. The state $\rho_{A'} \sigma$ satisfies
$v_j \wedge v_{n+j}$ for no $j$: for $j = i$ because $v_i = 0$ under
$\rho_{A'}$, and for $j \neq i$ because $v_{n+j} = 0$ under $\sigma$.
No other operator is applicable in either state, since every $o_i$ and
$\bar o_i$ requires $\neg x_i$ and $x_1, \dots, x_{2n}$ are all true.
Hence $\rho_{A'} \sigma$ is a dead
end and $\hstar(\rho_{A'} \sigma) = \infty \neq 1 =
\hstar(\rho_A \sigma)$, so
$\hstar|_{\rho_A} \neq \hstar|_{\rho_{A'}}$, and
$|\sub{m}{\hstar}| \geq 2^{n}$. Since the argument evaluates \hstar\
only on states, it applies to every total extension of \hstar\ to
the full cube $\{0,1\}^{4n+2}$ (Remark~\ref{rem-encoding}).
\end{proof}

\section{Width Bounds for Heuristic Families}

Let a SAS$^+$ task have variables $\vars=\{v_1,\dots,v_m\}$, let
$d=\max_v|D_v|$, and use a variable-contiguous Boolean order $\pi$.
An integer potential assigns $P(v,w)\in\mathbb Z$ to each fact, defines
$p(s)=\sum_{v\in\vars}P(v,s(v))$, and exposes the rectified heuristic
$\hpot(s)=\max(0,p(s))$.

\subsection{Integer Potentials}

\begin{proposition}
\label{prop-pot}
Let $p$ have integer potentials with $|P(v,w)|\leq M$ for all facts,
extended to non-state assignments by the same sum over the bit
encoding's induced table, with table values on invalid bit codes also
in $[-M,M]$, e.g., $0$, and let $\pi$ be variable-contiguous. Then
\[
\width_\pi(\hpot)\leq\width_\pi(p)\leq d(2mM+1)
\qquad\text{and}\qquad
|\img(\hpot)|\leq mM+1.
\]
\end{proposition}

\begin{proof}
Consider a prefix $\rho$ that fixes the variables
$v_1,\dots,v_j$ completely and the first
$b<\lceil\log_2d\rceil$ bits of $v_{j+1}$. The cofactor $p|_\rho$ is
determined by the accumulated sum
$\sum_{j'\leq j}P(v_{j'},\rho(v_{j'}))$, an integer in $[-mM,mM]$,
and by the assignment to the $b$ read bits of $v_{j+1}$. There are at
most $2mM+1$ sums and at most $2^b\leq d$ bit prefixes, so
$\width_\pi(p)\leq d(2mM+1)$. The raw total lies in $[-mM,mM]$;
rectification maps it to $\{0,\dots,mM\}$ and, as a terminal mapping,
cannot create cofactors.
\end{proof}

\begin{corollary}
\label{cor-pot}
For rectified consistent integer potential heuristics with magnitudes at
most $M$,
\[
\effort(\task,\hpot)
 \leq 3d(mM+1)(2mM+1)n\cdot\effort(\task,\hblind),
\]
a factor polynomial in the task size and the numeric magnitude $M$.
\end{corollary}

\begin{proof}
Proposition~\ref{prop-pot} gives
$W\leq d(2mM+1)$ and $\values\leq mM+1$; substitute both bounds in
Corollary~\ref{cor-effort}.
\end{proof}

\subsection{Pattern Databases}

\begin{proposition}
\label{prop-pdb}
Let $P\subseteq\vars$ be any pattern under a variable-contiguous order
$\pi$. Extend its table to invalid encodings by ignoring non-pattern
bits and mapping an invalid pattern-variable code to a fixed sink
terminal. Let
$|\states^P|=\prod_{v\in P}|D_v|$. Then
\[
\width_\pi(h^P)\leq d\,|\states^P|.
\]
\end{proposition}

\begin{proof}
Fix a bit-prefix $\rho$. Values read from variables outside $P$ are
irrelevant to this extension. If an already read pattern code is
invalid, the cofactor is the sink constant. Otherwise it is determined
by the assignments to the fully read pattern variables and, if the
current variable belongs to $P$, by its already-read bits. The number
of valid partial abstract assignments and bit buffers is at most
$|\states^P|$; adding the sink case still fits
$d|\states^P|$ for $d\geq2$ (the degenerate $d=1$ encoding has no
invalid bit code). Prefixes producing the same case have identical
completions. Once every pattern variable has been read, the cofactor is
a constant PDB value or the fixed sink, again within the stated bound.
\end{proof}

\subsection{Linear Merge-and-Shrink}

A linear merge-and-shrink mapping is a cascade of tables
$\tau_1,\dots,\tau_m$, where $\tau_j$ maps the previous intermediate
abstract state and the value of $v_j$ to the next intermediate state.

\begin{proposition}
\label{prop-ms}
Let $h^{\alpha}$ be an M\&S heuristic built with a linear merge
strategy whose variable order agrees with the variable-contiguous
$\pi$ and with intermediate abstraction sizes at most $N$. Extend each
table with a distinguished sink for invalid variable codes. Then
$\width_\pi(h^{\alpha}) \leq d \cdot (N+1)$.
\end{proposition}

\begin{proof}
For a prefix $\rho$ fixing $v_1, \dots, v_j$ and $b$ bits of
$v_{j+1}$, the cofactor $h^{\alpha}|_\rho$ is determined by the
intermediate abstract state $\tau_j(\cdots)$ reached after processing
$v_1, \dots, v_j$ through the cascading tables, of which there are at
most $N$, or the distinguished sink, together with the $b$-bit buffer,
at most $d$ cases: the
remaining computation applies the fixed tables
$\tau_{j+1}, \dots, \tau_m$ and the final distance table to the suffix.
Thus at every cut there are at most $d(N+1)$ distinct cofactors.
\end{proof}

\section{Width-Constrained Selection: Complete Implementation}
\label{sec-supp-selection}

The implementation treats width as a feasibility constraint and heuristic
information as the ranking objective. Under the total-ADD definition and
goal-zero convention, width one forces the single attained value zero, so the
$K=1$ selector deliberately includes the blind endpoint as a sanity check.
Two heuristics of the same larger width can still differ substantially in
initial estimate, dead-end recognition, construction time, and the BDDs
encountered during search. The design therefore combines exact measurement, a
fixed PDB candidate pool, same-$g$ batching, and a fail-safe construction
budget.

\subsection{Exact Width of a Total Heuristic ADD}

Every candidate is represented as one total ADD in the search manager's
variable order. Finite PDB distances are its ordinary terminals; abstract
dead ends, when present, receive a distinct sentinel terminal. Thus the
measurement covers the implementation's total extension rather than only its
finite level sets; bit assignments outside a valid abstract-state encoding
retain the default zero terminal.

To compute cofactor width exactly, we maintain a frontier of distinct residual
ADD nodes. The frontier initially contains the root. At each manager level, a
residual rooted at that level is replaced by its two children, while a
terminal or a residual whose root is below that level is retained. CUDD
unique-table identity is equality of residual functions, so deduplicating node
pointers after each step gives the exact cofactor count at that cut; the
maximum over the initial cut and all manager levels is $\width_\pi(h)$.
Retaining residuals at skipped levels is necessary because a reduced ADD need
not display a node at every variable. As cross-checks, the last frontier must
equal the set of distinct terminals and the measured width must not exceed the
total ADD node bound of Proposition~\ref{prop-add}. Dynamic reordering is
disabled: otherwise a width measured during candidate selection would not
certify the order used later by search.

\subsection{Fixed, Deduplicated PDB Pool and Score}

Given an abstract-state budget $B=100{,}000$, the selector constructs the same
five-source pool on every run: (i) the empty pattern; (ii) the longest prefix
of the BDD variable order that fits $B$; (iii) the longest prefix of a
goal/causal variable order that fits $B$; (iv) a \emph{goal-fill} pattern that
follows the same goal/causal order but skips a nonfitting variable and
continues with later ones; and (v) one pattern from Fast Downward's seeded
CEGAR generator \citep{rovner-et-al-icaps2019}. The CEGAR call uses seed 2011,
a 10-second process-time construction limit, and the same state budget. These
are candidate generators, not novel PDB algorithms.

Patterns are sorted, duplicate variables are removed, and identical patterns
from different generators are materialized only once while retaining all
source labels. This provenance-preserving deduplication is important at small
budgets, where several generators can collapse to the same pattern. Each
unique candidate within $B$ is materialized as an explicit PDB and a total
ADD, and is feasible exactly when its measured cofactor width is at most the
declared budget $K$. The empty pattern guarantees a feasible width-one
candidate. The selector never grows or repairs a candidate after observing
its width; it filters this fixed pool, which keeps comparisons across values
of $K$ well-defined.

Among feasible candidates, a deterministic lexicographic score prefers:
recognition of the initial state as a dead end; otherwise larger
$h(\sinit)$; larger mean finite abstract goal distance; larger abstract
dead-end fraction; smaller exact width; fewer abstract states; and finally
the lexicographically smaller sorted pattern. The first four components are
informativeness surrogates, while width is a feasibility constraint and only
a late tie-breaker. The scoring rule is declared in advance; it is not a
claim that these surrogates optimize runtime.

\subsection{Same-$g$ Speculative Batching}

Product-at-evaluation \bddastar\ normally invokes the instrumented high-level
cost-image routine
once for each selected $(g,h)$ bucket. With batching window $\Delta>0$, the
implementation forms an exact union of the selected bucket and eligible fresh
buckets having the \emph{same} $g$ and $f\leq F+\Delta$, where $F$ is the
selected minimum $f$, and computes one cost-image of that union. Restricting the
union to one $g$ is essential: the BDD would otherwise erase the parent
path-cost label needed to assign successor keys. Optional gates reject an
additional bucket if the candidate union exceeds an inner-node limit or a
union-to-separate-node ratio.

Batching is speculative only. Extra source buckets remain pending at their
original keys; their goal tests, insertion into the closed list, and logical
selection occur later in the ordinary $(f,g)$ order. The already-imaged
\emph{source} BDDs are parked until that selection, while successors of the
union image are partitioned and inserted immediately. For a consistent
heuristic and positive operator costs, the implementation invariant is that
those successors cannot overtake their pending sources. An independent
regression oracle compares the exact logical expansion trace and solution cost
with the unbatched path; batching may nevertheless generate successors that
unbatched search would never need. The case $\Delta=0$ is the ordinary
per-bucket path. The evaluation includes a fixed $\Delta=16$ treatment and
adaptive treatments with $\Delta=64$, maximum union ratio $1.0$, and maximum
union size $10^6$ inner nodes. Width bounds alone do not predict whether the
reduction in calls repays union and speculation costs.

\subsection{Construction Budgets and Safe Fallback}

M\&S exposes a separate setup risk: a bounded-width abstraction can still be
too expensive to build. The \texttt{build\_time\_limit} covers both the
linear merge-and-shrink construction and conversion of its final abstraction
to symbolic level sets. If the deadline expires or construction leaves more
than one active factor, the incomplete heuristic is discarded and the same
wrapper continues with blind forward symbolic search; construction time
remains charged to the run. A finite-value ceiling $\phi_\kappa\circ h$ is also
available and, by Proposition~\ref{prop-transform}, preserves consistency and
cannot increase width while reducing the finite value count.

The screening controls use a 10,000-state causal-graph linear abstraction with
\texttt{align\_merge\_order=false}, both with uncapped abstract goal distances
and with finite heuristic values capped at 32. Exact width is still measured
in the search order, but the closed-form $d(N+1)$ guarantee does not apply to
these controls. An aligned construction would have the a priori $d(N+1)$
bound but is a different abstraction-quality treatment. Integer-potential box
constraints remain another analytic route; they are outside the present
\texttt{release\_no\_lp} evaluation.

\section{Empirical Protocol and Frozen Provenance}
\label{sec-supp-protocol}

All reported results come from the current Arrhenius cluster; no older-cluster
run is pooled with or used as a substitute for this evaluation. Every
treatment uses the same forward symbolic framework and task-specific fixed BDD
variable order. Completed nonblind heuristics use the product-at-evaluation
\bddastar\ engine and minimum-$(f,g)$ bucket ordering; the blind control and
safeguard fallbacks use the equivalent $h\equiv0$ forward uniform-cost path.
Treatments otherwise differ only in the heuristic and in the explicit batching
or construction safeguard named below. Throughout, micro-PAR2 is the
arithmetic mean of reported planner CPU time for solved cells and 600 seconds
for every unsolved cell. Heuristic construction is inside the 300-second
planner limit and is therefore charged to that time.

\subsection{Questions and Compared Methods}

We predeclare three questions. \textbf{Q1:} How does the exact width budget
change the selected PDB's width, pattern size and source, and search outcomes?
\textbf{Q2:} Does same-$g$ batching reduce completed high-level cost-image
calls and image time without sacrificing coverage? \textbf{Q3:} Does the resulting
selector improve on conventional goal-fill and CEGAR PDB construction, and how
does it compare with blind and M\&S controls after construction time is charged?
The study is diagnostic: the completed P4/P5 evidence is negative for the
practical usefulness of the current selector and batching designs, while the
pending P6 census remains unresolved and is not prejudged here.

The screening matrix contains five controls and fifteen selector treatments.
The controls are blind forward search; causal-graph linear M\&S with a
10,000-state shrink limit and uncapped abstract goal distances; the same M\&S
configuration with finite heuristic values capped at 32; goal-fill PDB
construction with $B=100{,}000$; and CEGAR PDB construction with the same $B$,
10-second limit, and seed 2011. Both M\&S controls set
\texttt{align\_merge\_order=false}; their measured exact width is informative,
but the $d(N+1)$ bound does not apply.

Ten selector treatments use
\[
 K\in\{1,2,4,8,16,32,64,128,256,\infty\}
\]
without batching. Four more cross $K\in\{32,64,128,\infty\}$ with the adaptive
batching treatment $\Delta=64$, union ratio at most $1.0$, and union size at
most $10^6$ inner nodes. The final treatment uses $K=32$ and the fixed window
$\Delta=16$ without the adaptive gates. Every selector treatment uses the
identical five-source, $B=100{,}000$ candidate pool and CEGAR settings described
above. The pool trace, its provenance-preserving deduplication, every candidate
score and width, the feasible set, the selected record, and deterministic
trace and pool hashes are logged. Pruning-only search is not part of this
frozen empirical matrix and remains an analytic design direction.

\subsection{Population and Stages}

The source inventory contains 1697 tasks from the optimal sequential IPC
benchmarks through 2023. Frozen, disjoint exclusions remove 290 tasks with
zero-cost operators and 30 with normalized axioms, leaving the target
population of 1377 positive-cost, axiom-free tasks. Positive operator costs
match the assumptions of the bucket-identity lemma and the batching argument;
the axiom exclusion is an implementation restriction of the abstraction
heuristics. The frozen support manifest establishes operator costs with 1684
direct SAS-v3 scans and, for 13 resource-heavy translations, a pinned-parser
no-metric/unit-cost proof. A separate attestation checks axioms after the
pinned translator's default normalization. Every selected PDDL source is
byte-attested, and tasks outside this population are rejected before launch.

\paragraph{Stage 1: frozen P4 screen.}
The corrected screening manifest contains 50 tasks from 25 domains, exactly
two tasks per domain. Crossing it with all 20 configurations gives 1000 cells.
Its composition was assembled during development using archived outcome
summaries to enrich difficult or discriminating cases. Those old measurements
influence the screen composition only: they are neither pooled with the
Arrhenius runs nor used as performance observations. The frozen ordering
records both the 20-configuration screen winner and the first selector-family
treatment in that same ordering. The latter remains available for a direct
test of width-constrained selection even if a control wins among the screen
configurations. This stage performs selection; it is not used as a
confirmatory estimate of either winner's advantage.

\paragraph{Stage 2: frozen P5 disjoint validation sample.}
A separately frozen 92-task validation manifest contains two supported tasks
in each of 46 eligible domains and is disjoint from the P4 screen,
cap-validation, and development-smoke tasks. Within each domain, eligible
remaining problem names are ordered by SHA-256 of a fixed seed, the domain,
and the problem name
(separated by NUL bytes), and the first two are selected; no planner outcome
enters this ordering. Six configurations are fixed independently of the
screen: the five screening controls plus M\&S capped at finite heuristic value
32 with a 60-second construction budget and blind fallback. The Stage-1
selector-family winner is unbatched and is added without retuning, so no
additional same-$K$ unbatched ablation is required. The screen winner is
already one of the fixed controls and is recorded as an alias rather than a
duplicate cell. The frozen stage therefore has exactly seven unique
configurations and 644 runs.

Here \emph{held out} means held out from P4 selection. We do not claim that
these validation tasks were historically unseen over the broader project.

The primary contrast is the selector-family winner against the CEGAR PDB
control. Its estimand is the paired coverage difference within each domain,
averaged equally over the 46 domains; a predeclared deterministic
100,000-replicate domain bootstrap supplies a 95\% domain-resampling interval.
The construction budget has its own matched safeguard. The bootstrap resamples
these 46 fixed domains. Its interval describes variation in domain
composition; it does not capture task-selection or run-to-run uncertainty and
is not an interval for the 1377-task census. Comparisons with blind search,
uncapped and capped M\&S, task-micro coverage, PAR2, planner CPU time,
construction time, image time, and observed exact-width summaries are
descriptive because there is one run per cell. No validation outcome changes
the configuration ordering. The screen artifact and validation rerun intentionally
retain separate provenance pins: the screen is the immutable selection record,
whereas the validation run uses the prospectively corrected producer revision
described below. We therefore compare treatments only within a stage and do
not interpret differences in absolute counts or times between stages as a
cross-stage effect.

\paragraph{Stage 3: frozen, pending P6 full census.}
The population, matrix, execution contract, acceptance rule, and P6 launch
identity are frozen, while all Stage-3 results remain pending. The complete
seven-configuration validation matrix is carried forward unchanged to all 1377
supported tasks on the same cluster, giving 9639 cells. A separate runner
freezes the selection-artifact hash, task-source bytes, option matrix, and
array layout. This stage is an exact descriptive census of the frozen
supported population, not a third sample-based validation stage. It
intentionally includes the 50 screening and 92 validation tasks. We report the
selector-versus-CEGAR coverage difference averaged equally over the 46
domains, overall and per-domain coverage, and operational cost without
sampling intervals. One predeclared sensitivity repeats that macro-coverage
contrast on the 1327 tasks outside the selection screen; we do not describe
this complement as historically unseen. No Stage-2 or Stage-3 outcome reopens
the selector, width budget, batching parameters, matrix, ranking, or contrasts.
Thus the screen performs selection, the disjoint validation set evaluates the
frozen choice, and the full census describes its scope.

\subsection{Resources, Staging, and Acceptance}

Every run has a 300-second aggregate process-CPU planner limit, 8\,GiB memory
limit, and one AMD EPYC 9755 core on Arrhenius. We use the
\texttt{release\_no\_lp} build, so the present matrix does not evaluate
potential heuristics. The scheduler permits at most five simultaneous array
elements. The screen and validation stages place one run in each element. To
stay below Arrhenius's 1000-element array limit, the frozen full-suite layout
deterministically executes at most 10 runs per element in 964 array elements,
with a 70-minute scheduler envelope; runs within an element remain sequential.
Time spent constructing a heuristic is inside the planner limit and therefore
charged to that configuration. This envelope is a conservative request, not a
per-cell wall-clock guarantee: the wrapper has no separate wall watchdog, and
an outer scheduler timeout invalidates the census.

After grid construction, every stage replaces each benchmark symlink with an
independent read-only copy, verifies each copy against the frozen per-source
digest, and reattests all cells immediately before submission; later changes
to the benchmark worktree therefore cannot change queued inputs. The
full-suite job additionally requests scheduler requeue for node failure or
preemption. On restart it skips only cells bearing an exact, atomically written
wrapper-completion marker and archives all dynamic state of an unmarked cell
before trying that cell again. The marker records a zero exit from the Python
run wrapper; it does not certify that the planner solved the task. The census
is accepted only if every declared cell is fetched exactly once with a valid
marker and a protocol-recognized planner outcome. A missing, starved,
duplicate, or unexpected infrastructure cell invalidates the census rather
than changing a denominator or becoming a timeout. Any recovery beyond
automatic scheduler requeue requires a separately frozen recovery artifact
before outcomes are inspected; there is no post-hoc cell fill.

\subsection{Frozen Selection and Measurement}

The Stage-1 promotion rule is one lexicographic order over all 20 labels:
first maximize solved count; then minimize suite-micro PAR2, charging every
unsolved cell 600 seconds; then minimize total image time only if all 1000
cells have complete, semantically certified schema-v2 counters; finally choose
the lexicographically smaller label. The screen winner is the first label in
this order; the selector-family winner is the first selector treatment in the
same order, without a second ranking rule. Costs of every solved task must
agree across configurations. Unexpected exit outcomes, missing cells, a
changed option string, or a provenance mismatch invalidates the screen rather
than silently becoming a timeout. Resource-killed runs may contribute
certified prefix counters, but an incomplete counter never enters the
image-time tie-break.

Schema v2 counts inner BDD nodes per stored piece, distinguishes attempted
from completed expansions and image calls, records source-bucket and
source-piece multiplicities, and measures construction and image time. For
each heuristic it also records finite value count, distinct terminals, total
ADD inner nodes, the $A+T$ upper bound, exact cofactor counts at every manager
level, and exact width. Selector logs are checked as a whole trace: source
order, deduplication, candidate feasibility, score recomputation, selected
winner, and final heuristic hashes must agree. This prevents a favorable
summary from being accepted if the pool or selection process drifted.
Selection and description are emitted separately. The version-3 selection
artifact carries the frozen ranking and downstream matrix. A canonical
\texttt{arrhenius-selector-screen-certified-descriptive-census/v2} report
reuses those exact public selection objects and adds certified pattern-size
and source counts plus image and batching totals only over explicitly counted
complete cells. Adaptive-gate acceptance and rejection summaries remain only
in console logs, so we do not report them as fetched protocol data.

The driver records both the raw search-process exit and the effective planner
exit. A raw resource-limit exit is promoted to a ``plan found plus resource
limit'' outcome only when that invocation created exactly one stable,
single-link, canonical plan. The independent parser then requires the exact
raw-to-effective mapping, agreement of plan-file, run-log, and metric-stream
costs, and a complete solved metric summary. Thus a CPU-limit signal during
successful teardown cannot silently turn a durable plan into either an
ordinary timeout or an uncertified success.

\subsection{P4/P5/P6 Identities and Prospective Corrections}

The reconciliation was frozen before the accepted screen. Two earlier
Arrhenius diagnostic grids are excluded: the first revealed that two zero-cost
tasks had passed a source-identity-only support gate; the second exposed one
durable plan followed by a raw CPU-limit exit during teardown. We
prospectively corrected task admission and raw/effective reconciliation,
rebuilt the 1000-cell grid, and reran every cell as P4. Those diagnostic
outcomes serve only as protocol-failure evidence; no score or performance
observation from them enters the selection artifact.

The source and execution identities are frozen independently. P4 supplies the
screen, P5 the accepted validation run, and P6 the full rerun; the corresponding
planner/launch revisions begin \texttt{58a3f742d}/\texttt{d6d98ad42},
\texttt{165b6d2ee}/\texttt{a52488637}, and
\texttt{a3486a027}/\texttt{0cb19e111}. Every runner byte-attests its planner,
preprocessor, code, PDDL inputs, run scripts, static metadata, and complete job
script before submission. The analyzers independently pin those identities,
the task and option matrices, and the canonical selection artifact. The
accompanying anonymous artifact contains the raw SHA-256 sidecars, exact
revision and protocol strings, runners, analyzers, immutable JSON results, and
a machine-readable rendering path; \texttt{paper/reproducibility-manifest.md}
indexes them in one place.

Two downstream audits caused prospective whole-grid corrections. First, a
heuristic that proved the initial state dead could exit before emitting its
completed construction event. Both arrays were canceled (53 stable validation
cells and 20 full-suite completion markers), the producer ordering was fixed,
and both grids were rebuilt from scratch; the resulting P5 validation grid is
the accepted one. Second, four of 1650 completion-marked P5 full cells received
\texttt{SIGXCPU} in legacy preprocessing, whose negative child return surfaced
through the shell as 232. P6 maps this pre-search condition to canonical
resource outcome 21 and again reruns all 9639 cells. A scheduled compute-node
smoke exercised that path before launch. No completed cell, partial cell, or
archived attempt from either rejected start is reused.

These amendments were triggered by protocol violations, not performance
comparisons. They leave the selection, seven-configuration matrix, ranking,
contrasts, and statistical policy unchanged. P6 records every declared
outcome-code bucket, including zeros; the unit is one retained, attested final
cell record, not a shell attempt generation or component-process invocation.
The full census remains pending.

\section{Complete Preliminary Results}
\label{sec-supp-results}

\subsection{P4 Screen}

All 1000 declared cells passed the protocol checks.
Table~\ref{tab:supp-screen} shows the decisive rows of the frozen 20-by-50
census. Uncapped M\&S and its control capped at finite heuristic value 32 both
solved 43 tasks, but uncapped M\&S had 1.3706 seconds lower micro-PAR2 and
therefore won the frozen screen ordering before the image-time tie-break.
Against blind search it gained 13 tasks net, with 17 exclusive wins and four
losses. All 950 heuristic constructions completed, with no construction or
pattern fallback.

\begin{table}[t]
  \centering
  \small
  \caption{Selected screening-census rows. Micro-PAR2 is in seconds with a
  600-second penalty. Width gives minimum/median/maximum exact cofactor width
  over the 50 tasks. ``Image $n$'' is the number of raw-complete, certified
  cells admitted to that configuration's total-image summary.}
  \label{tab:supp-screen}
  \begin{tabular}{lrrcr}
    \toprule
    Configuration & Solved & Micro-PAR2 & Width & Image $n$ \\
    \midrule
    Blind forward & 30 & 264.9006 & -- & 30 \\
    Uncapped M\&S (screen winner) & 43 & 133.6736 & $11/201/149{,}079$ & 43 \\
    M\&S, finite-value cap 32 & 43 & 135.0442 & $11/185.5/58{,}010$ & 43 \\
    Selector $K=1$ (family winner) & 30 & 273.2554 & $1/1/1$ & 30 \\
    CEGAR PDB & 27 & 312.1382 & $2/18/2{,}696$ & 27 \\
    \bottomrule
  \end{tabular}
\end{table}

The selector-family winner, $K=1$, tied blind coverage but increased
micro-PAR2 by 8.3548 seconds; descriptively, it solved three more tasks than
the CEGAR PDB control. Its 50 certified traces all selected the empty,
size-zero pattern and had exact width one. This PDB is semantically
$h\equiv0$: it induces the blind partition while retaining the heuristic
wrapper's construction and evaluation overhead. Relaxing the unbatched budget
did not improve this result: the nine $K>1$ treatments solved 28--30 tasks and
had micro-PAR2 277.1670--291.7778. At $K=32$, for example, selected widths
ranged from 1 to 30 (median 17), selected pattern sizes from 0 to 16, and the
primary-source counts in pool order---empty, BDD prefix, goal prefix, goal
fill, and CEGAR---were $8/12/18/1/11$. Thus larger safe widths changed the
selected abstractions without improving the selector-family ordering. Across
the ten unbatched budgets in increasing order, the number of empty-pattern
selections fell from 50 to $47/40/31/21/8/5/1/0/0$, yet coverage never exceeded
30 tasks.

Adaptive batching also regressed on this screen. Each $\Delta=64$ treatment
lost three tasks against its same-$K$ unbatched reference and increased
micro-PAR2 by 33.8674--44.8040 seconds; the fixed $\Delta=16$, $K=32$
treatment lost six tasks and increased micro-PAR2 by 57.4396 seconds. On the
raw-complete certified intersections, batching reduced completed image calls
to 11.8--23.2\% of the corresponding unbatched counts, but image time was
1.31--2.06 times as large. These comparisons contain only 24--27 of the 50
task pairs; the other 23--26 pairs are censored.

Every cell nevertheless supplied a piece-certified prefix, and all 750
selector traces were complete. Only 589 of 1000 cells were raw-complete for
total-image aggregation, however, leaving 411 censored; consequently the
ranking's total image times are undefined and the image-time tie-break was not
used. These are descriptive selection-screen results only and make no claim
about the validation sample or the full supported population.

\subsection{P5 Disjoint Validation Sample}

All 644 validation cells passed the protocol checks. In the primary comparison,
the width-one selector solved 41 tasks and CEGAR solved 40. Because every
domain contributes exactly two tasks, the equally weighted domain-macro and
task-micro estimands coincide in this sample; the selector-minus-CEGAR
difference under either estimand is $1/92=1.09$ percentage points. The
predeclared 100,000-replicate domain-bootstrap 95\% domain-resampling interval
was $[-1/46,1/23]$, or $[-2.17,4.35]$ percentage points, with two task-level
discordant wins and one loss. Thus the validation data do not establish a
selector coverage advantage over CEGAR. The selector's micro-PAR2 was 345.1040
seconds against 351.5588 seconds for CEGAR, while its planner CPU-time
geometric-mean ratio on the 39 jointly solved tasks was 1.006. Because all
eligible traces again selected the empty PDB, this primary contrast evaluates
the frozen selection procedure in its blind-equivalent $h\equiv0$ outcome; it
does not validate selection of an informative heuristic.

\begin{table}[t]
  \centering
  \small
  \caption{P5 validation results on 92 tasks held out from P4 selection.
  PAR2 charges 600 seconds to an
  unsolved cell. Width reports observed-cell count followed by median/maximum
  exact width. Construction reports observed records/fallbacks and their total
  time. Gaps are a pre-search cell in every configuration and, for uncapped
  M\&S, one additional resource prefix before construction; the five
  build-budget fallbacks supply no M\&S width.}
  \label{tab:supp-heldout}
  \begin{tabular}{lrrcrr}
    \toprule
    Configuration & Solved & Micro-PAR2 & Width $n;\widetilde{x}/\max$ &
      Constr. $n/f$ & Constr. s \\
    \midrule
    Blind forward & 40 & 348.1627 & -- & -- & -- \\
    Uncapped M\&S & 42 & 338.7962 & $90;142/14{,}540$ & $90/0$ & 769.928 \\
    M\&S, finite-value cap 32 & 42 & 336.8272 & $91;135/14{,}540$ & $91/0$ & 1017.722 \\
    Goal-fill PDB & 41 & 345.3886 & $91;31/6{,}670$ & $91/0$ & 27.074 \\
    CEGAR PDB & 40 & 351.5588 & $91;17/6{,}672$ & $91/0$ & 12.570 \\
    M\&S, value cap 32, 60-s build & 43 & 334.1384 & $86;123/14{,}540$ & $91/5$ & 701.311 \\
    Selector $K=1$ & 41 & 345.1040 & $91;1/1$ & $91/0$ & 54.185 \\
    \bottomrule
  \end{tabular}
\end{table}

The construction safeguard had the highest coverage among the seven evaluated
forward configurations in Table~\ref{tab:supp-heldout}: it solved 43 tasks,
one more than capped M\&S without a build deadline (one discordant win, no
loss), and reduced micro-PAR2 by 2.6888 seconds. Five of its 92 cells fell back
to blind search (5.43\%). Across the 91 paired construction observations its
total construction time was 0.689 times the reference total; on the 42 jointly
solved tasks its planner CPU-time geometric-mean ratio was 0.994. On the 86
cells where both M\&S constructions completed, their paired exact-width totals
were identical (90,613); on the 42 jointly solved cells, their BDD-node-effort
totals were also identical. These aggregate equalities are consistent with
preserving the capped abstraction when the deadline is inactive. The five
fallbacks identify cells on which construction did depart, but the frozen
aggregate artifact does not attribute the discordant solve to a particular
fallback. These are diagnostics of the predeclared safeguard, not a new
ranking.

Every configuration contributed 91 semantically certified image prefixes and
one pre-search-empty cell. Complete-summary image counts in table order were
$41/43/43/42/41/44/42$, with corresponding total image times
$477.883/341.866/321.073/549.118/483.450/572.440/673.078$ seconds. On the
primary 40-pair complete-summary intersection, however, the selector used
350.098 seconds against CEGAR's 446.812 seconds, a ratio of 0.784. The same two
methods had 91 paired construction and width observations: selector
construction took 4.311 times the CEGAR total, whereas its paired width totals
were 91 and 14,707. On 39 jointly solved cells its BDD-node effort ratio was
0.679. Thus the selector spent 4.311 times CEGAR's paired construction total
to return $h\equiv0$ in every eligible trace. All 91 eligible selector traces
were complete and certified.

For context, the selector tied goal-fill coverage (two discordant wins and two
losses) and lowered PAR2 by 0.2846 seconds; against blind it gained one task
(one win, no loss) and lowered PAR2 by 3.0587 seconds; against uncapped and
capped M\&S it lost one task (two wins and three losses in each comparison)
and increased PAR2 by 6.3078 and 8.2768 seconds. The respective jointly solved
planner CPU-time geometric-mean ratios were 0.874 ($n=39$), 1.428 ($n=40$),
0.908 ($n=39$), and 0.907 ($n=39$). Uncapped M\&S itself gained two tasks over
blind (four wins and two losses) and lowered PAR2 by 9.3665 seconds, although
its jointly solved planner CPU-time ratio was 1.632 ($n=38$). The seven frozen
configurations and every declared comparison remain unchanged; no P5
validation outcome was used to rerank them.

Taken together, the completed P4/P5 results are a negative diagnostic for the
current selector and batching designs: the frozen selector repeatedly returns
the blind-equivalent empty PDB without an established CEGAR advantage, while
the tested batching variants lose coverage and raise image time. This neither
validates usefulness nor predicts the pending P6 census.

\subsection{P6 Full-Suite Result Placeholder}

\emph{[RESULT PLACEHOLDER: in the main paper, use one combined P5/P6 table
with the same seven frozen rows rather than a separate P5 table. Report the
current-Arrhenius descriptive census on all 1377 supported tasks, including
the selector-versus-CEGAR equally weighted
46-domain macro-coverage contrast; overall and per-domain coverage; all-task
micro-PAR2; solved planner CPU-time summaries; and the carried jointly solved
or paired-observation comparisons. Report certified exact-width and
construction summaries, including capped-M\&S fallback incidence;
certified-prefix image time; complete-certified image time and call counts;
and complete-call amortization; complete-trace selector effort and selector
selected-source counts; and solved-cell BDD-node effort, each with its exact
eligible, certified/observed, solved, complete-summary, or paired-intersection
denominator. Report the sole predeclared screen-excluded 1327-task sensitivity
without an interval, plus the protocol-recognized retained final-cell outcome
census, completion-marker validity, Slurm-restart, partial-archive, and shell
cell-attempt-generation counts globally and per configuration. Distinguish
final cell records from attempt generations, and do not describe either as
planner invocations. Do not rerank configurations or substitute any archived
coverage table.]}

Until this placeholder is filled from the attested P6 artifact, the
population-level empirical conclusion remains deliberately unresolved. The
formal safety results do not depend on any empirical outcome.

\section{Additional Related-Work Detail}

Heuristic symbolic search has long exposed a tension between information and
representation. \bddastar\ partitions a frontier by absolute heuristic value
\citep{edelkamp-reffel-ki1998}; SetA\ensuremath{^{*}} and \ghseta\ move
heuristic changes into transition-relation partitions
\citep{jensen-et-al-aij2008}; and ADDA\ensuremath{^{*}} represents values
algebraically \citep{hansen-et-al-sara2002}. The exponential separation of
\citet{speck-et-al-icaps2020} shows that even \hstar\ can be harmful under the
BDD-node effort measure. Our response is deliberately narrower than a general
performance theorem: cofactor width is a sufficient representation-safety
condition for product-at-evaluation forward search. The exact-width selector
operationalizes that condition, but neither its ranking score nor its runtime
is covered by the theorem.

PDB construction and symbolic PDB use are established techniques
\citep{culberson-schaeffer-compint1998,edelkamp-ecp2001}. In particular,
automatic CEGAR pattern generation is due to
\citet{rovner-et-al-icaps2019}, and compiling PDBs and M\&S heuristics into
BDD value sets predates this work \citep{edelkamp-et-al-ecai2012}.
\citet{torralba-phd2015} studied polynomial BDD representations of linear M\&S
heuristics, while \citet{helmert-et-al-icaps2015} made their cascading
factored mappings explicit. We do not claim the empty, BDD-prefix,
goal/causal-prefix, goal-fill, or CEGAR generators as new heuristic
algorithms. The methodological contribution is the fixed,
provenance-preserving deduplicated pool, exact measurement of each resulting
total ADD in the search order, and rejection before ranking whenever the
declared width budget is exceeded.

Potential heuristics \citep{pommerening-et-al-aaai2015} and their integration
into symbolic search \citep{fiser-et-al-aaai2022,fiser-et-al-aij2024} provide
another compact family. Corollary~\ref{cor-pot} bounds rectified integer
potential level sets, not the cost-and-delta transition partitions used by
\ghseta. Likewise, Proposition~\ref{prop-ms} turns the M\&S shrink limit into
the bound $d(N+1)$ only when the linear merge order agrees with the search
variable order. A causal-graph construction with
\texttt{align\_merge\_order=false} is not certified by that formula, even if
its width is measured exactly afterward.

Aggregating BDD work is not new. SetA\ensuremath{^{*}} groups compatible state
sets in a best-set-first queue \citep{jensen-et-al-aij2008}; List
\bddastar\ stores a whole $g$-row in one BDD and defers extraction of its
heuristic slices \citep{torralba-et-al-aij2017}; and transition trees merge
same-cost transition relations subject to a BDD-size limit
\citep{torralba-et-al-icaps2013}. Our same-$g$ batching makes a narrower trade:
already separated source buckets with different $h$-values are united only for
one speculative image, while every source remains pending at its original key.
It then replays the original logical $(f,g)$ order. This addresses the
image-call multiplier of Proposition~\ref{prop-calls} without changing the
heuristic representation, permanently merging open-list buckets, or compiling
heuristic deltas into transition relations. Trace and solution-cost
preservation under consistency and positive costs are implementation
invariants checked by an exact regression oracle; batching carries no a priori
runtime guarantee.

Finally, \citet{speck-helmert-ecai2025} compare symbolic and explicit blind
search through the representation of transition relations and goal formulas.
Our comparison is complementary: it bounds heuristic over blind symbolic
representation effort. Variable-ordering studies
\citep{kissmann-hoffmann-jair2014} are directly relevant because width is
order-dependent. Symbolic perimeter abstractions
\citep{torralba-et-al-ijcai2016a,torralba-et-al-aij2018} and pruning-only
search both delay or restrict heuristic work, but use different mechanisms.
Theorem~\ref{thm-bucket} applies to a fixed slice in either search direction;
a cumulative bidirectional guarantee accounting for changing incumbents,
meeting points, and direction choices remains open.

\bibliographystyle{plainnat}
\bibliography{bib/abbrv,bib/literatur,bib/crossref}
\end{document}
